\centerline{\bf \Huge Многочлены}

\begin{flushright}
        \scriptsize{я хочу спать}
\end{flushright}

\problem{1}
В выражении $\left(x^4-2 x^3+3 x^2-4 x+5\right)^{2016}$ раскрыли скобки и привели подобные слагаемые. Найдите сумму коэффициентов получившегося многочлена.

\problem{2}
$a$ и $b$ - положительные числа. Сумма минимального значения квадратного трехчлена $f(x)=a x^2+8 x+b$ и минимального значения квадратного трехчлена $g(x)=b x^2+8 x+a$ равна нулю. Докажите, что эти минимальные значения оба равны нулю.

\problem{3}
Приведенные квадратные трехчлены $f(x)$ и $g(x)$ принимают неположительные значения на непересекающихся отрезках. Докажите, что найдутся такие положительные числа $\alpha$ и $\beta$, что для любого действительного $x$ будет выполняться неравенство $\alpha f(x)+\beta g(x)>0$.

\problem{4}
Приведенный квадратный трехчлен с целыми коэффициентами в трех последовательных целых точек принимает простые значения. Докажите, что он принимает простое значение по крайней мере еще в одной целой точке.

\problem{5}
Дан квадратный трехчлен $2015 x^2+2016 x+2017$. Двое ходят по очереди. Своим ходом игрок может вычесть из многочлена $x^2, x$ или 1 . Проигрывает игрок, после хода которого получается многочлен с целочисленным корнем. Может ли один из игроков обеспечить себе победу?

\problem{6}
Существуют ли такие четыре квадратных трехчлена, что, в каком бы порядке $f_1, f_2, f_3, f_4$ их ни выписали, найдется такое вещественное число $a$, что $f_1(a)<f_2(a)<f_3(a)<f_4(a)$ ?

\problem{7}
Докажите, что при любых натуральных $k, l, m$ многочлен $x^{3 k+2}+x^{3 l+1}+ x^{3 m}$ делится на $x^2+x+1$.

\problem{8}
Даны приведённые многочлены $P(x)$ и $Q(x)$ степени 2016. Известно, что уравнение $P(x)=Q(x)$ не имеет действительных корней. Докажите, что уравнение $P(x+1)=Q(x-1)$ имеет хотя бы один действительный корень.

\problem{9}
Пусть $P(x)$ - многочлен нечетной степени, имеющий 2017 различных корней. Докажите, что уравнение $P(P(x))=0$ также имеет хотя бы 2017 различных корней.